\chapter{The ele_struct}
\label{c:ele.struct}
\index{ele_struct|hyperbf}

This chapter describes the \vn{ele_struct} which is the structure that
holds all the information about an individual lattice element:
quadrupoles, separators, wigglers, etc. The \vn{ele_struct} structure is
shown in \figs{f:ele.struct1}  and \ref{f:ele.struct2}. This
structure is somewhat complicated, however, in practice, a lot of the
complexity is generally hidden  by the \bmad bookkeeping routines.

As a general rule, for variables like the Twiss parameters that are not
constant along the length of an element, the value stored in the
corresponding component in the \vn{ele_struct} is the value at the downstream
end of the element.

For printing information about an element, the
\Hyperref{r:type.ele}{type_ele} routine can be used (\sref{s:first.program}).
When called without the optional \vn{lines} and \vn{n_lines} arguments, \vn{type_ele}
prints to the terminal window; when those arguments are supplied, the element
information is returned as an array of strings instead.
\nopagebreak[4]
\begin{figure}[htb]
\centering
\footnotesize
\begin{verbatim}
type ele_struct
  character(40) name                       ! name of element \sref{c:ele.string}.
  character(40) type                       ! type name \sref{c:ele.string}.
  character(40) alias                      ! Another name \sref{c:ele.string}.
  character(40) component_name             ! Used by overlays, multipass patch, etc.
  character(200), pointer :: descrip       ! Description string.
  type (twiss_struct)  a, b, z             ! Twiss parameters at end of element \sref{c:normal.modes}.
  type (xy_disp_struct) x, y               ! Projected dispersions \sref{c:normal.modes}.
  type (ac_kicker_struct), pointer :: ac_kick  ! ac_kicker element parameters.
  type (bookkeeping_state_struct) bookkeeping_state ! Element attribute bookkeeping
  type (branch_struct), pointer :: branch      ! Pointer to branch containing element.
  type (controller_struct), pointer :: control ! For group, overlay, and ramper elements.
  type (converter_struct), pointer :: converter  ! Positron converter, etc.
  type (rf_ele_struct), pointer :: rf            ! RF element stair-step parameters.
  type (foil_struct), pointer :: foil            ! Foil material parameters.
  type (ele_struct), pointer :: lord           ! Pointer to a slice lord.
  type (fibre), pointer :: ptc_fibre           ! PTC tracking.
  type (floor_position_struct) floor           ! Global floor position.
  type (high_energy_space_charge_struct), pointer :: high_energy_space_charge
                                               ! High-energy space charge parameters.
  type (mode3_struct), pointer :: mode3        ! Full 6-dimensional normal mode decomposition.
  type (multipole_cache_struct), allocatable :: multipole_cache  ! Cached multipole values.
  type (photon_element_struct), pointer :: photon
  type (rad_map_ele_struct), pointer :: rad_map ! Radiation damping and stochastic kick maps.
  type (taylor_struct) :: taylor(6)            ! Orbital Taylor map.
  real(rp) :: spin_taylor_ref_orb_in(6)        ! Reference orbit for spin Taylor map.
  type (taylor_struct) :: spin_taylor(0:3)     ! Spin Taylor map (quaternion representation).
  type (wake_struct), pointer :: wake          ! Wakes.
  type (wall3d_struct), pointer :: wall3d(:)   ! Chamber or capillary wall array.
  ! E/M field map components (see Chapter on Electromagnetic Fields):
  type (cartesian_map_struct), pointer :: cartesian_map(:)     ! Used to define E/M fields.
  type (cylindrical_map_struct), pointer :: cylindrical_map(:) ! Used to define E/M fields.
  type (gen_grad_map_struct), pointer :: gen_grad_map(:)       ! Used to define E/M fields.
  type (grid_field_struct), pointer :: grid_field(:)           ! Used to define DC and AC fields.
    ele_struct definition continued on next figure...
\end{verbatim}
\caption[The \vn{ele_struct} (part 1).]{The \vn{ele_struct}. structure definition. 
The complete structure is shown in this and the following figure.}
\label{f:ele.struct1}
\end{figure}

%--------------------------------------------------------------------------

\begin{figure}[tb]
\centering
\footnotesize
\begin{verbatim}
    ... ele_struct definition continued from previous figure.
  type(coord_struct) map_ref_orb_in            ! Transfer map ref orbit at upstream end of element.
  type(coord_struct) map_ref_orb_out           ! Transfer map ref orbit at downstream end of element.
  type(coord_struct) time_ref_orb_in           ! Reference orbit at upstream end for ref_time calc.
  type(coord_struct) time_ref_orb_out          ! Reference orbit at downstream end for ref_time calc.
  real(rp) value(num_ele_attrib$)              ! attribute values.
  real(rp) old_value(num_ele_attrib$)          ! Used to see if %value(:) array has changed.
  real(rp) spin_q(0:3,0:6)                     ! 0th and 1st order spin transport quaternion.
  real(rp) vec0(6)                             ! 0th order transport vector.
  real(rp) mat6(6,6)                           ! 1st order transport matrix.
  real(rp) c_mat(2,2)                          ! 2x2 C coupling matrix
  real(rp) dc_mat_dpz(2,2)                     ! d(c_mat)/dpz chromatic variation.
  real(rp) gamma_c                             ! gamma associated with C matrix
  real(rp) s_start                             ! longitudinal ref position at entrance_end
  real(rp) s                                   ! longitudinal position at the downstream end.
  real(rp) ref_time                            ! Time ref particle passes downstream end.
  real(rp), pointer :: r(:,:,:)                ! For general use. Not used by Bmad.
  real(rp), pointer :: a_pole(:)               ! multipole a_n (skew) components.
  real(rp), pointer :: b_pole(:)               ! multipole b_n (normal) components.
  real(rp), pointer :: a_pole_elec(:)          ! Electrostatic multipoles.
  real(rp), pointer :: b_pole_elec(:)          ! Electrostatic multipoles.
  real(rp), pointer :: custom(:)               ! Custom attributes
  integer key                    ! key value
  integer sub_key                ! Records bend input type (rbend$, sbend$).
  integer ix_ele                 ! Index in lat%branch(n)%ele(:) array [n = 0 <==> lat%ele(:)].
  integer ix_branch              ! Index in lat%branch(:) array [0 => In lat%ele(:)].
  integer lord_status            ! overlay_lord$, etc.
  integer n_slave                ! Number of slaves (not counting field overlap slaves).
  integer n_slave_field          ! Number of field slaves.
  integer ix1_slave              ! Pointer to lat%control array.
  integer slave_status           ! super_slave$, etc.
  integer n_lord                 ! Number of lords (not counting field overlap and ramper lords).
  integer n_lord_field           ! Number of field lords.
  integer n_lord_ramper          ! Number of ramper lords.
  integer ic1_lord               ! Pointer to lat%ic array.
  integer ix_pointer             ! For general use. Not used by Bmad.
  integer ixx, iyy, izz          ! Index for Bmad internal use.
  integer mat6_calc_method       ! bmad_standard$, taylor$, etc.
  integer tracking_method        ! bmad_standard$, taylor$, etc.
  integer spin_tracking_method   ! bmad_standard$, symp_lie_ptc$, etc.
  integer csr_method             ! off$, one_dim$, steady_state_3d$, etc.
  integer space_charge_method    ! off$, slice$, fft_3d$, cathode_fft_3d$, etc.
  integer ptc_integration_type   ! drift_kick$, matrix_kick$, etc.
  integer field_calc             ! bmad_standard$, no_field$, fieldmap$, custom$, etc.
  integer ref_species            ! Reference species (not_set$ by default).
  integer aperture_at            ! Aperture location: exit_end$, entrance_end$, etc.
  integer aperture_type          ! Type of aperture: rectangular$, elliptical$, or custom$.
  integer orientation            ! -1 -> Element is longitudinally reversed. +1 -> Normal.
  logical symplectify            ! Symplectify mat6 matrices.
  logical mode_flip              ! Have the normal modes traded places?
  logical multipoles_on          ! For turning multipoles on/off.
  logical scale_multipoles       ! Multipole components scaled by strength of element?
  logical taylor_map_includes_offsets  ! Taylor map calculated with element misalignments?
  logical field_master           ! Calculate strength from the field value?
  logical is_on                  ! For turning element on/off.
  logical logic                  ! For general use. Not used by Bmad.
  logical bmad_logic             ! For Bmad internal use only.
  logical select                 ! For element selection. Used by make_hybrid_ring, etc.
  logical offset_moves_aperture  ! Element offsets affect aperture position?
end type
\end{verbatim}
\caption[The \vn{ele_struct} (part 2).]{The \vn{ele_struct}. 
The complete structure is shown in this and the preceding figure.}
\label{f:ele.struct2}
\end{figure}

%--------------------------------------------------------------------------
\section{Initialization and Pointers}
\index{ele_struct!initialization}
\index{ele_struct!pointer components}

The \vn{ele_struct} has a number of components and subcomponents that are pointers. This is done to
minimize memory size. The use of pointers raises a deallocation issue.  Generally, most \vn{ele_struct}
elements are part of a \vn{lat_struct} variable (\sref{s:lat:point}) and such elements in a
\vn{lat_struct} are handled by the \vn{lat_struct} allocation/deallocation routines.  In the case
where a local \vn{ele_struct} variable is used within a subroutine or function, the \vn{ele_struct}
variable must either be defined with the \vn{save} attribute
\begin{example}
  type (ele_struct), save :: ele          ! Use the save attribute
  logical, save :: init_needed = .false.
  ...
  if (init_needed) then
    call init_ele (ele, quadrupole\$)     ! Initialize element once
    init_needed = .false.
  endif
\end{example}
or the pointers within the variable must be deallocated  with a call to
\Hyperref{r:deallocate.ele.pointers}{deallocate_ele_pointers}:
\begin{example}
  type (ele_struct) ele  
  ...
  call init_ele (ele, sbend\$)            ! Initialize element each time
  ...
  call deallocate_ele_pointers (ele)     ! And deallocate.
\end{example}

In the ``normal'' course of events, the pointers of an \vn{ele_struct} variable should not be
pointing to the same memory locations as the pointers of any other \vn{ele_struct} variable. To make
sure of this, the equal sign in the assignment \vn{ele1 = ele2} is overloaded by the routine
\Hyperref{r:ele.equal.ele}{ele_equal_ele}.  The exception here are the ``Electro-magnetic field
component'' pointer arrays \vn{ele%cartesian_map}, \vn{ele%cylindrical_map},
\vn{ele%gen_grad_map}, and \vn{ele%grid_field}. Since these components potentially contain large
arrays, and since the individual sub-components are not likely to be individually modified, the
field component pointers of \vn{ele1} and \vn{ele2} after the set \vn{ele1 = ele2} will point at
the same memory locations.

Note: The assignment \vn{ele1 = ele2} will not modify \vn{ele1%ix_ele} or \vn{ele1%ix_branch}. If
\vn{ele1} is associated with a lattice then \vn{ele1%lat} will also be unaffected.

%--------------------------------------------------------------------------
\section{Miscellaneous}
\label{s:ele.misc}

The \vn{%ref_species} integer component stores the reference particle species for an element
This can differ from element to element in a branch due to \vn{converter} or \vn{foil} elements.
Note: The setting of \vn{branch%param%particle} gives the reference species at the beginning of
a branch and this will correspond to the setting of \vn{%ref_species} of the beginning element
of that branch.

%--------------------------------------------------------------------------
\section{Element Attribute Bookkeeping}
\label{s:ele.book}
\index{Element attribute bookkeeping}

When a value of an attribute in an element changes, the values of other attributes may need to be
changed (\sref{s:depend}). Furthermore, in a lattice, changes to one element may necessitate changes
to attribute values in other elements. For example, changing the accelerating gradient in an
\vn{lcavity} will change the reference energy throughout the lattice.

The attribute bookkeeping for a lattice can be complicated and, if not done intelligently, can cause
programs to be slow if attributes are continually being changed. In order to keep track what
bookkeeping has been done, the \vn{ele%status} component is used by the appropriate bookkeeping
routines for making sure the bookkeeping overhead is keep to a minimum. See \sref{s:lat.book} for
more details.

%--------------------------------------------------------------------------
\section{String Components}
\label{s:ele.string}

\index{ele_struct!\%descrip}\index{ele_struct!\%alias}
\index{ele_struct!\%type}\index{ele_struct!\%name}
The \vn{%name}, \vn{%type}, \vn{%alias}, and \vn{%descrip} components of the \vn{ele_struct} all
have a direct correspondence with the \vn{name}, \vn{type}, \vn{alias}, and \vn{descrip} element
attributes in an input lattice file (\sref{s:alias}). On input (\sref{s:lat.readin}), The \vn{name}
attribute is converted to uppercase before being loaded into an \vn{ele_struct}. The other three are
not. To save memory, since \vn{%descrip} is not frequently used, \vn{%descrip} is a pointer that is
only allocated if \vn{descrip} is set for a given element.

When a lattice is constructed by \Hyperref{r:bmad.parser}{bmad_parser}, a ``nametable'' is
constructed to enable the fast lookup of element names. Therefore, It is important that if element
names are modified, that the nametable is updated appropriately. This is done by calling
\Hyperref{r:create.lat.ele.nametable}{create_lat_ele_nametable} after any element name
modifications.

%--------------------------------------------------------------------------
\section{Element Key}
\label{s:ele.key}

\index{ele_struct!\%key}
The \vn{%key} integer component gives the class of element
(\vn{quadrupole}, \vn{rfcavity}, etc.). In general, to get the
corresponding integer parameter for an element class, just add a ``\$''
character to the class name. For example \vn{quadrupole\$} is the integer
parameter for \vn{quadrupole} elements. The \vn{key_name} array converts from
integer to the appropriate string. For example:
\begin{example}
  type (ele_struct) ele
  if (ele%key == wiggler\$) then       ! Test if element is a wiggler.
  print *, "This element: ", key_name(ele%key) ! Prints, for example, "WIGGLER"
\end{example}
Note: The call to \vn{init_ele} is needed for any \vn{ele_struct} defined
outside of a \vn{lat_struct} structure.

\index{wiggler}\index{rbend}\index{sbend}\index{ele_struct!\%sub_key}
The \vn{%sub_key} component is only used for bend element.  When a lattice file is parsed,
(\sref{s:lat.readin}), all \vn{rbend} elements are converted into \vn{sbend} elements
(\sref{s:bend}). To keep track of what the original definition of the element was, the \vn{%sub_key}
component will be set to \vn{sbend\$} or \vn{rbend\$} whatever is appropriate. 
The \vn{%sub_key} component does not affect any calculations and is only used in the routines that
recreate lattice files from a \vn{lat_struct} (\sref{s:lat.write}).

%--------------------------------------------------------------------------
\section{The \%value(:) array}
\label{s:ele.dep}
\index{ele_struct!attribute values}

\index{ele_struct!\%value(:)}
Most of the real valued attributes of an element are held in the \vn{%value(:)} array. For example,
the value of the \vn{k1} attribute for a quadrupole element is stored in \vn{%value(k1\$)} where
\vn{k1\$} is an integer parameter that \bmad defines.  In general, to get the correct index in
\vn{%value(:)} for a given attribute, add a ``\$" as a suffix. To convert from an attribute name to
its index in the \vn{%value} array use the \Hyperref{r:attribute.index}{attribute_index} routine.
To go back from an index in the \vn{%value} array to a name use the
\Hyperref{r:attribute.name}{attribute_name} routine. Example:
\begin{example}
  type (ele_struct) ele
  call init_ele (ele, quadrupole\$)    ! Initialize element
  ele%value(k1\$) = 0.3                                        ! Set K1 value
  print *, "Index for Quad K1:  ", attribute_index(ele, "K1") ! prints: `4' (= k1\$)
  print *, "Name for Quad k1\$: ", attribute_name (ele, k1\$)   ! prints: `K1' 
\end{example}
The list of attributes for a given element
type is given in the writeup for the different element in
Chapter~\ref{c:elements}. 

To obtain a list of attribute names and associated \vn{%value(:)}
indexes, the program \vn{element_attributes} can be used. This program
is included in the standard \bmad distribution.

Besides real valued attributes, the \vn{value(:)} array also holds logical,
integer, and, as explained below, ``switch'' attributes. To find out the
type of a given attribute, use the function \Hyperref{r:attribute.type}{attribute_type}. See the
routine \vn{type_ele} for an example of how \vn{attribute_type} is used.

An example of a logical attribute is the \vn{flexible} logical of match elements which is stored in
\vn{%value(flexible\$)}. To evaluate logical attributes, the functions
\Hyperref{r:is.true}{is_true(param)} or \Hyperref{r:is.false}{is_false(param)} should be used.

Integer attributes stored in the \vn{value(:)} array include \vn{n_slice} (stored
in \vn{%value(n_slice\$)}). With integer attributes, the \vn{nint(param)} 
Fortran instrinsic should be used for evaluation.

A switch attribte is an attribute whose value is one of a certain set of integers
where each integer corresponds to some ``state''. For example, the \vn{fringe_at}
switch which, as explained in \sref{s:fringe}, may have one of four values.
Generally, the integer parameters that correspond to the states of a switch
can be constructed by putting a ``\$'' after the associated name. Thus,
with the \vn{fringe_at} switch, the four integer parameters are \vn{no_end\$},
\vn{both_ends\$}, \vn{entrance_end\$}, and \vn{exit_end\$}. For example:
\begin{example}
  if (nint(ele%value(fringe_type$)) == soft_edge_only$) then
    ...
\end{example}
For printing purposes, to convert a switch value to the appropriate string, use
the routine\Hyperref{r:switch.attrib.value.name}{switch_attrib_value_name} can be used.

\index{ele_struct!\%field_master}\index{harmon}
The \vn{%field_master} logical within an element sets whether it is the normalized strength or field
strength that is the independent variable. See \sref{s:depend} for more details.

\index{ele_struct!\%old_value(:)}
The \vn{%old_value(:)} component of the \vn{ele_struct} is used by the
\Hyperref{r:attribute.bookkeeper}{attribute_bookkeeper} routine to check for changes for changes in
the \vn{%value(:)} array since the last time the \vn{attribute_bookkeeper} routine had been called.
If nothing has been changed, the \vn{attribute_bookkeeper} routine knows not to waste time
recalculating dependent values. Essentially what this means is that the \vn{%old_value(:)} array
should not be modified outside of \vn{attribute_bookkeeper}.

%--------------------------------------------------------------------------
\section{Connection with the Lat_Struct}
\label{s:ele.lat}

\index{ele_struct!\%ix_ele}
\index{ele_struct!\%ix_branch}
\index{ele_struct!\%branch}
If an element is part of a \vn{lat_struct} (\sref{c:lat.struct}), the
\vn{%ix_ele} and \vn{%ix_branch} components of the \vn{ele_struct}
identify where the element is. Additionally, the \vn{%branch} component
points to the \vn{branch_struct} containing the element, and the
enclosing \vn{lat_struct} is accessed via \vn{ele%branch%lat}. That is
\begin{example}
  type (lat_struct), pointer :: lat
  type (ele_struct), pointer :: ele2
  if (ele%ix_ele > -1) then
    ie = ele%ix_ele
    ib = ele%ix_branch
    lat => ele%branch%lat
    ele2 => lat%branch(ib)%ele(ie)
    print *, associated(ele2, ele)  ! Will print True.
  endif
\end{example}
In this example the \vn{ele2} pointer is constructed to point to the \vn{ele} element. The test
(\vn{ele%ix_ele}~>~-1) is needed since \vn{ele_struct} elements may exist outside of any
\vn{lat_struct} instance. Such ``external'' elements always have \vn{%ix_ele}~<~0.  A value for
\vn{%ix_ele} of \vn{ix_slice_slave\$} (= -2) is special in that it prevents the
\Hyperref{r:deallocate.ele.pointers}{deallocate_ele_pointers} routine from deallocating the pointers
of an element which has its \vn{%ix_ele} set to \vn{ix_slice_slave\$}.

An element ``slice'' is an example of an element that exists external to any \vn{lat_struct}
instance. A slice is an \vn{ele_struct} instance that represents some sub-section of a given
element. Element slices are useful when tracking particles only part way through an element
(\sref{s:tracking.partial}).

%--------------------------------------------------------------------------
\section{Limits}
\index{aperture}
\index{x_limit}\index{x1_limit}\index{x2_limit}
\index{x_limit}\index{x1_limit}\index{x2_limit}

The aperture limits (\sref{s:limit}) in the \vn{ele_struct} are:
\begin{example}
  %value(x1_limit$)
  %value(x2_limit$)
  %value(y1_limit$)
  %value(y2_limit$)
\end{example}
The values of these limits along with the \vn{%aperture_at}, \vn{%aperture_type},
and \vn{%offset_moves_aperture} components are used in tracking to determine
if a particle has hit the vacuum chamber wall. See Section~\sref{s:tracking.apertures}
for more details.

%--------------------------------------------------------------------------
\section{Twiss Parameters, etc.}
\label{s:ele.twiss}

\index{ele_struct!\%gamma_c}
\index{ele_struct!\%c_mat}
\index{ele_struct!\%a}\index{ele_struct!\%b}\index{ele_struct!\%z}
\index{ele_struct!\%mode3}\index{ele_struct!\%mode_flip}

The components \vn{%a}, \vn{%b}, \vn{%z}, \vn{%x}, \vn{%y}, \vn{%c_mat},
\vn{%gamma_c}, \vn{%mode_flip}, and \vn{%mode3}
hold information on the Twiss parameters, dispersion, and coupling
at the downstream end of the element. See \cref{c:normal.modes} for more details.

\index{ele_struct!\%dc_mat_dpz}
The \vn{%dc_mat_dpz(2,2)} component stores the chromatic derivative of the $2 \times 2$
C-coupling matrix: $\partial C / \partial p_z$. This is computed along with the C matrix
during the normal mode analysis.

%--------------------------------------------------------------------------
\section{Element Lords and Element Slaves}
\label{s:ele.control}

\index{ele_struct!\%slave_status}
\index{ele_struct!\%n_slave}
\index{ele_struct!\%n_slave_field}
\index{ele_struct!\%ix1_slave}
\index{ele_struct!\%lord_status}
\index{ele_struct!\%n_lord}
\index{ele_struct!\%n_lord_field}
\index{ele_struct!\%n_lord_ramper}
\index{ele_struct!\%ic1_lord}
\index{ele_struct!\%component_name}
In \bmad, elements in a lattice can control other elements.
The components that determine this control are:
\begin{example}
  %slave_status
  %n_slave          ! Number of slaves (not counting field overlap slaves)
  %n_slave_field    ! Number of field overlap slaves
  %ix1_slave
  %lord_status
  %n_lord           ! Number of lords (not counting field overlap and ramper lords)
  %n_lord_field     ! Number of field overlap lords
  %n_lord_ramper    ! Number of ramper lords
  %ic1_lord
  %component_name
\end{example}
\index{lat_struct}
This is explained fully in the chapter on the \vn{lat_struct} (\sref{c:lat.struct}).

%--------------------------------------------------------------------------
\section{Group and Overlay Controller Elements}
\label{s:control.ele}
\index{group}\index{overlay}

\vn{Group}, \vn{overlay}, and \vn{ramper} elements use the \vn{%control} pointer for storing
controller information. \vn{%control} points to a \vn{controller_struct} which holds the control
variables and ramp expressions:
\begin{example}
  type controller_struct
    type (control_var1_struct), allocatable :: var(:)      ! Control variables
    type (control_ramp1_struct), allocatable :: ramp(:)    ! For ramper lord elements
    type (ramper_lord_struct), allocatable :: ramper_lord(:) ! Ramper lord info for this slave
    real(rp), allocatable :: x_knot(:)
  end type
\end{example}
The \vn{var(:)} component is an array of \vn{control_var1_struct} elements, one per control variable:
\begin{example}
  type control_var1_struct
    character(40) :: name = ""
    real(rp) :: value = 0
    real(rp) :: old_value = 0
  end type
\end{example}
The \vn{%old_value} component of \vn{control_var1_struct} is only used for \vn{group} elements.

See Section~\sref{s:lat.add.delete} for an example of setting up a controller element within a program.

%--------------------------------------------------------------------------
\section{Coordinates, Offsets, etc.}
\label{s:ele.offset}
\index{global coordinates!in ele_struct}

The ``upstream'' and ``downstream'' ends of an element are, by definition,
where the physical ends of the element would be if there were no
offsets. In particular, if an element has a finite \vn{z_offset}, the
physical ends will be displaced from upstream and downstream ends. See
\sref{s:track1} for more details.

\index{ele_struct!\%floor}
\index{floor_position_struct}
The \vn{%floor} component gives the ``laboratory'' global ``floor'' coordinates (\sref{s:global})
at the downstream end of the element. These coordinates are computed without misalignments.
That is, the coordinates are not ``body'' coordinates. The components of the \vn{%floor} structure are
\begin{example}
  type floor_position_struct
    real(rp) r(3)               ! Offset from origin
    real(rp) w(3,3)             ! Orientation matrix (\Eq{wwww})
    real(rp) theta, phi, psi    ! Angular orientation
  end type
\end{example}
The routine \Hyperref{r:ele.geometry}{ele_geometry} will calculate an element's 
floor coordinates given the floor coordinates at the beginning of the element.
In a lattice, the \Hyperref{r:lat.geometry}{lat_geometry} routine will calculate
the floor coordinates for the entire lattice using repeated calls to \vn{ele_geometry}.

\index{x_offset}\index{y_offset}\index{x_pitch}\index{y_pitch}
\index{tilt}\index{x_offset_tot}\index{y_offset_tot}\index{x_pitch_tot}
\index{y_pitch_tot}\index{tilt_tot}
The positional offsets (\sref{s:offset}) for an element 
from the reference orbit are stored in
\begin{example}
  %value(x_offset\$)
  %value(y_offset\$)
  %value(z_offset\$)
  %value(x_pitch\$)
  %value(y_pitch\$)
  %value(tilt\$)
\end{example}
\index{girder}
If the element is supported by a \vn{girder} element (\sref{s:girder})
then the \vn{girder} offsets are added to the element offsets 
and the total offset with respect to the
reference coordinate system is stored in:
\begin{example}
  %value(x_offset_tot\$)
  %value(y_offset_tot\$)
  %value(z_offset_tot\$)
  %value(x_pitch_tot\$)
  %value(y_pitch_tot\$)
  %value(tilt_tot\$)
\end{example}
If there is no \vn{girder}, the values for \vn{%value(x_offset_tot\$)}, etc.
are set to the corresponding values in \vn{%value(x_offset\$)}, etc.
Thus, to vary the position of an individual
element the values of \vn{%value(x_offset\$)}, etc. are changed and to
read the position of an element a program should look at
\vn{%value(x_offset_tot\$)}, etc.

\index{ele_struct!\%s}
\index{ele_struct!\%ref_time}
The longitudinal position at the downstream end of an element is stored in \vn{%s} and the reference
time is stored in \vn{%ref_time}. This reference time is calculated assuming that the reference time
is zero at the start of the lattice. Also stored in the \vn{ele_struct} is the reference time at the
start of the element and the differnece in the reference time between the end and the beginning.
These are given in \vn{%value(ref_time_start\$)} and \vn{%value(delta_ref_time\$)} respectively.

Notice that the reference time used to calculate the $z$ phase
space coordinate (\Eq{zbctt}) may be different from \vn{%ref_time}. For example, with multiple
bunches the $z$ phase space coordinate is generally taken to be with respect to a reference particle
at the center of the bunch the particle is in. And, at a given element, the reference time of the
different bunch reference particles will be different. Another example happens when a particle is
tracked through multiple turns. In this case the reference time at a given element will depend upon
the turn number.

%--------------------------------------------------------------------------
\section{Transfer Maps: Linear and Non-linear (Taylor)}
\label{s:ele.maps}
\index{ele_struct!transfer maps}
\index{transfer map!in ele_struct}
\index{ele_struct!\%mat6}
\index{ele_struct!\%vec0}
\index{ele_struct!\%map_ref_orb_in}
\index{ele_struct!\%map_ref_orb_out}

The routine \Hyperref{r:make.mat6}{make_mat6} computes the linear 
transfer matrix (Jacobian) along with the zeroth order transfer vector. 
This matrix is stored in \vn{%mat6(6,6)} and the
zeroth order vector is stored in \vn{%vec0(6)}. The reference orbit at
the upstream end of the element about
which the transfer matrix is computed is stored in \vn{%map_ref_orb_in}
and the reference orbit at the downstream end is stored in \vn{%map_ref_orb_out}.
In the calculation of the transfer map, the vector \vn{%vec0} is set so that
\begin{example}
  map_ref_orb_out = %mat6 * map_ref_orbit_in + %vec0
\end{example}
The reason redundant information is stored in the element is to save
computation time.

To compute the transfer maps for an entire lattice use the routine 
\Hyperref{r:lat.make.mat6}{lat_make_mat6}.

\index{ele_struct!Taylor maps}
\index{taylor map!structure in ele_struct}
The Taylor map (\sref{c:methods})  for an element is stored in
\vn{%taylor(1:6)}. Each \vn{%taylor(i)} is a \vn{taylor_struct}
structure that defines a Taylor series:
\begin{example}
  type taylor_struct
    real (rp) ref
    type (taylor_term_struct), pointer :: term(:) => null()
  end type
\end{example}
Each Taylor series has an array of \vn{taylor_term_struct} terms defined as
\begin{example}
  type taylor_term_struct
    real(rp) :: coef
    integer :: expn(6)
  end type
\end{example}
The coefficient for a Taylor term is stored in \vn{%coef} and the
six exponents are stored in \vn{%exp(6)}. 

To see if there is a Taylor map associated with an element the
association status of \vn{%taylor(1)%term} needs to be checked.
As an example the following finds the order of a Taylor map.
\begin{example}
  type (ele_struct) ele
  ...
  if (associated(ele%taylor(1)%term) then  ! Taylor map exists
    taylor_order = 0
    do i = 1, 6
      do j = 1, size(ele%taylor(i)%term)
        taylor_order = max(taylor_order, sum(ele%taylor(i)%term(j)%exp)
      enddo
    enddo
  else  ! Taylor map does not exist
    taylor_order = -1  ! flag non-existence
  endif
\end{example}

The Taylor map is made up around some reference phase space point
corresponding to the coordinates at the upstream of the element.
This reference point is saved in \vn{%taylor(1:6)%ref}.  Once a Taylor map is
made, the reference point is not needed in subsequent
calculations. However, the Taylor map itself will depend upon what
reference point is chosen (\sref{s:taylor.phys}).

%--------------------------------------------------------------------------
\section{Reference Energy and Time}

The reference energy and reference time are computed around a
reference orbit which is different from the reference orbit used for
computing transfer maps (\sref{s:ele.maps}). The energy and time
reference orbit for an element is stored in
\begin{example}
  ele%time_ref_orb_in        ! Reference orbit at upstream end
  ele%time_ref_orb_out       ! Reference orbit at downstream end
\end{example}
Generally \vn{ele%time_ref_orb_in} is the zero orbit. The exception
comes when an element is a \vn{super_slave}. In this case, the
reference orbit through the super_slaves of a given \vn{super_lord} is
constructed to be continuous. This is done for consistancey sake. For
example, to ensure that when a marker is superimposed on top of a
wiggler the reference orbit, and hence the reference time, is not altered.

%--------------------------------------------------------------------------
\section{EM Field Maps}
\label{s:ele.em.field}

\index{ele_struct!\%cartesian_map}
\index{ele_struct!\%cylindrical_map}
\index{ele_struct!\%gen_grad_map}
\index{ele_struct!\%grid_field}
The electromagnetic field of an element may be described using one or more field map
representations.  These are stored in four pointer array components of the \vn{ele_struct}:
\begin{example}
  ele%cartesian_map(:)     ! Cartesian harmonic field decomposition
  ele%cylindrical_map(:)   ! Cylindrical harmonic field decomposition
  ele%gen_grad_map(:)      ! Generalized gradient field model
  ele%grid_field(:)        ! Sampled-grid field (DC or AC)
\end{example}
Since each of these is a pointer array, its association status must be tested before accessing
any sub-components. Multiple field map types may be present simultaneously on a single element.
Full documentation of these structures and the underlying field models is given in
Chapter~\ref{c:electromagnetic.fields}.

%--------------------------------------------------------------------------
\section{Wakes}
\label{s:ele.wake}
\index{wakefields!in ele_struct}

\index{ele_struct!\%wake}
\index{lcavity}\index{rfcavity}
The \vn{ele%wake} component holds information on the wakes associated with an element.  Since
\vn{ele%wake} is a pointer, its association status must be tested before any of its sub--components
are accessed.
\begin{example}
  type (ele_struct) ele
  ...
  if (associated(ele%wake)) then
    ...
\end{example}
\bmad observes the following rule: If \vn{%wake} is associated, it is assumed that all the
sub--components (\vn{%wake%sr_table}, etc.) are associated. This simplifies programming in that you
do not have to test directly the association status of the sub--components.

\index{ele_struct!\%wake}
See \sref{c:wakefields.eq} for the equations used in wakefield calculations.  wakefields are stored
in the \vn{%wake} struct:
\begin{example}
  type wake_struct
    type (wake_sr_struct) :: sr     ! Short-range wake.
    type (wake_lr_struct) :: lr     ! Long-range wake.
    real(rp) :: amp_scale = 1       ! Wake amplitude scale factor.
    real(rp) :: time_scale = 1      ! Wake time scale factor.
  end type
\end{example}

\index{wake_sr_mode_struct}
The short--range wake parameterization uses pseudo--modes (\sref{s:sr.wake.eq}). This parameterization
utilizes the \vn{%wake%sr%long(:)}, and \vn{%wake%sr%trans(:)} arrays for the longitudinal and
transverse modes respectively.  The structure used for the elements of these arrays are:
\begin{example}
  type wake_sr_mode_struct  ! Pseudo-mode short-range wake struct 
    real(rp) amp        ! Amplitude
    real(rp) damp       ! Damping factor.
    real(rp) freq       ! Frequency in Hz
    real(rp) phi        ! Phase in radians/2pi
    real(rp) b_sin      ! non-skew sin-like component of the wake
    real(rp) b_cos      ! non-skew cos-like component of the wake
    real(rp) a_sin      ! skew sin-like component of the wake
    real(rp) a_cos      ! skew cos-like component of the wake
  end type
\end{example}
The wakefield kick is calculated from \Eq{wadzk}. \vn{%amp}, \vn{%damp}, \vn{%freq}, and \vn{%phi}
are the input parameters from the lattice file. the last four components (\vn{%norm_sin}, etc.)
store the accumulated wake: Before the bunch passes through these are set to zero and as each
particle passes through the cavity the contribution to the wake due to the particle is calculated
and added the components.

\vn{%wake%z_sr_mode_max} is the maximum $z$ value beyond which the pseudo mode representation is not
valid. This is set in the input lattice file.

\index{wake_lr_struct}
The \vn{%wake%lr} array stores the long--range wake modes. The
structure definition is:
\begin{example}
  type wake_lr_struct   ! Long-Range Wake struct 
    real(rp) freq       ! Actual Frequency in Hz
    real(rp) freq_in    ! Input frequency in Hz
    real(rp) R_over_Q   ! Strength in V/C/m^2
    real(rp) Q          ! Quality factor
    real(rp) angle      ! polarization angle (radians/2pi).
    integer m           ! Order (1 = dipole, 2 = quad, etc.)
    real(rp) b_sin      ! non-skew sin-like component of the wake
    real(rp) b_cos      ! non-skew cos-like component of the wake
    real(rp) a_sin      ! skew sin-like component of the wake
    real(rp) a_cos      ! skew cos-like component of the wake
    logical polarized   ! Polarized mode?
  end type
\end{example}
\index{lr_freq_spread}
This is similar to the \vn{sr_mode_wake_struct}. \vn{%freq_in} is the actual frequency in the input
file. \Hyperref{r:bmad.parser}{bmad_parser} will set \vn{%freq} to \vn{%freq_in} except when the
\vn{lr_freq_spread} attribute is non-zero in which case \vn{bmad_parser} will vary \vn{%freq} as
explained in \sref{s:lcav}. \vn{%polarized} is a logical that indicates whether the mode has a
polarization angle. If so, then \vn{%angle} is the polarization angle.

%--------------------------------------------------------------------------
\section{Wiggler and Undulator Fields}
\index{wiggler!field maps}

Wiggler and undulator elements (\vn{wiggler} and \vn{undulator} element types) define their
magnetic fields using the standard field map components of the \vn{ele_struct} (see
\sref{s:ele.em.field}): \vn{%cartesian_map}, \vn{%cylindrical_map}, \vn{%gen_grad_map}, or
\vn{%grid_field}. The field decomposition models are described in detail in
Chapter~\ref{c:electromagnetic.fields}.

%--------------------------------------------------------------------------
\section{Spin Transport}
\label{s:ele.spin}
\index{ele_struct!\%spin_q}
\index{ele_struct!\%spin_taylor}
\index{ele_struct!\%spin_taylor_ref_orb_in}

The linear spin transport through an element is stored in the \vn{%spin_q(0:3,0:6)} array, which
gives the zeroth and first order quaternion transfer map. The reference orbit about which
this map is computed is the same as \vn{%map_ref_orb_in}/\vn{%map_ref_orb_out}
(\sref{s:ele.maps}).

For higher-order spin transport, the full spin Taylor map is stored in
\vn{%spin_taylor(0:3)}, a quaternion representation as an array of four
\vn{taylor_struct} elements. The reference orbit for the spin Taylor map
may differ from the orbital Taylor map reference orbit; it is stored separately
in \vn{%spin_taylor_ref_orb_in(6)}.

%--------------------------------------------------------------------------
\section{Radiation Maps}
\label{s:ele.rad.map}
\index{ele_struct!\%rad_map}

The \vn{%rad_map} pointer stores radiation damping and stochastic kick matrices for an element.
It is of type \vn{rad_map_ele_struct}:
\begin{example}
  type rad_map_ele_struct
    type (rad_map_struct) rm0, rm1  ! Upstream and downstream half matrices.
    logical :: stale = .true.
  end type
\end{example}
Each \vn{rad_map_struct} contains:
\begin{example}
  type rad_map_struct
    real(rp) :: ref_orb(6)           ! Reference orbit for the map.
    real(rp) :: damp_dmat(6,6)       ! Damping correction matrix.
    real(rp) :: xfer_damp_vec(6)     ! 0th order transfer vector with damping.
    real(rp) :: xfer_damp_mat(6,6)   ! 1st order matrix with damping.
    real(rp) :: stoc_mat(6,6)        ! Stochastic (Cholesky-decomposed) kick matrix.
  end type
\end{example}
Since \vn{ele%rad_map} is a pointer, its association status must be tested before use.

%--------------------------------------------------------------------------
\section{High-Energy Space Charge}
\label{s:ele.space.charge}
\index{ele_struct!\%high_energy_space_charge}

The \vn{%high_energy_space_charge} pointer stores the beam parameters needed for high-energy
space charge calculations. It is of type \vn{high_energy_space_charge_struct}:
\begin{example}
  type high_energy_space_charge_struct
    type (coord_struct) closed_orb   ! Beam orbit
    real(rp) kick_const
    real(rp) sig_x, sig_y, sig_z
    real(rp) phi                     ! Rotation angle, lab to rotated frame.
    real(rp) sin_phi, cos_phi
  end type
\end{example}
Since \vn{ele%high_energy_space_charge} is a pointer, its association status must be tested
before accessing sub-components.

%--------------------------------------------------------------------------
\section{Converter, Foil, and RF Components}
\label{s:ele.converter.foil.rf}
\index{ele_struct!\%converter}
\index{ele_struct!\%foil}
\index{ele_struct!\%rf}

Three pointer components of the \vn{ele_struct} are used by specialized element types:

\vn{%converter} is used by elements that convert one particle species to another (for example,
a positron converter target in a linac). It points to a \vn{converter_struct} holding output
species information and a table of particle distributions as a function of incoming energy and
material thickness.

\vn{%foil} is used by \vn{foil} elements and points to a \vn{foil_struct} describing the
material composition of the foil.

\vn{%rf} is used by RF cavity elements to store the stair-step energy model (\vn{rf_ele_struct}).
This model divides the cavity into a sequence of energy kicks interspersed with drifts. Note that
\vn{ele%rf} is not allocated for slice and super slaves.

Since all three are pointers, their association status must be tested before use.

%--------------------------------------------------------------------------
\section {Multipoles}
\index{ele_struct!multipoles}

\index{ele_struct!\%b_pole(:)}\index{ele_struct!\%a_pole(:)}
\index{multipole!an, bn!in ele_struct}
\index{multipole!KnL, Tn!in ele_struct}
\index{multipole!\%scale_multipoles}
The multipole components of an element (See \sref{s:mag.field}) are
stored in the pointers \vn{%a_pole(:)} and \vn{%b_pole(:)}. If
\vn{%a_pole} and \vn{%b_pole} are allocated they always have a range
\vn{%a_pole(0:n_pole_maxx)} and \vn{%b_pole(0:n_pole_maxx)}. Currently
\vn{n_pole_maxx} = 20. For a \vn{Multipole} element, the
\vn{%a_pole(n)} array stores the integrated multipole strength
\vn{KnL}, and the \vn{%b_pole(n)} array stores the tilt \vn{Tn}.

A list of \bmad routines for manipulating multipoles can be found in
\sref{r:multipoles}.

%--------------------------------------------------------------------------
\section{Tracking Methods}
\label{s:ele.tracking.methods}

A number of \vn{ele_struct} components control tracking and transfer
map calculations.  These are:
\begin{example}
  %mat6_calc_method           ! Integer switch: bmad_standard$, taylor$, etc.
  %tracking_method            ! Integer switch: bmad_standard$, taylor$, linear$, etc.
  %spin_tracking_method       ! Integer switch: bmad_standard$, symp_lie_ptc$, etc.
  %field_calc                 ! Integer switch: bmad_standard$, no_field$, fieldmap$, etc.
  %csr_method                 ! Integer switch: off$, one_dim$, steady_state_3d$, etc.
  %space_charge_method        ! Integer switch: off$, slice$, fft_3d$, cathode_fft_3d$, etc.
  %symplectify                ! Logical: symplectify mat6 after calculation?
  %multipoles_on              ! Logical: include multipole components?
  %taylor_map_includes_offsets ! Logical: Taylor map includes element misalignments?
  %is_on                      ! Logical: element on or off?
  %offset_moves_aperture      ! Logical: does element offset move the aperture?
\end{example}
Note that \vn{%csr_method} and \vn{%space_charge_method} are \vn{integer} switch components
(not logicals). See Chapter~\sref{c:tracking} for more details.

%--------------------------------------------------------------------------
\section{Custom and General Use Attributes}
\index{ele_struct!components not used by Bmad}
\label{s:ele.gen}

\index{ele_struct!\%r}
\index{ele_struct!\%ix_pointer}
\index{ele_struct!\%logic}
There are four components of an \vn{ele_struct} that are guaranteed to never be used by any \bmad
routine and so are available for use by someone writing a program. These components are:
\begin{example}
  %r(:,:,:)                     ! real(rp), pointer.
  %custom(:)                    ! real(rp), pointer.
  %ix_pointer                   ! integer.
  %logic                        ! logical.
\end{example}
Values for \vn{ele%r} and \vn{ele%custom} can be set in the lattice file (\sref{s:cust.att}). If
values are set for \vn{ele%r} or \vn{ele%custom}, these arrays will be expanded in size if needed.

Accessing the \vn{ele%custom} array should be done using the standard accessor routines
(\sref{s:lat.ele.attribute}). For example, if the lattice file being used defines a custom
attribute called \vn{rise_time} (\sref{s:cust.att}):
\begin{example}
  parameter[custom_attribute1] = "rise_time"
\end{example}
then a program can access the \vn{rise_time} attribute via:
\begin{example}
  type (ele_struct), poiner :: ele
  type (all_pointer_struct) a_ptr
  ...
  ele => ...            ! Point ele to some element in the lattice
  call pointer_to_attribute (ele, "RISE_TIME", .true., a_ptr, err_flag)
  print *, "RISE_TIME Attribute has value:", a_ptr%r
\end{example}

Note: Even if there are custom attributes associated with a given type of element (say, all
quadrupoles), a given element of that type may not have its \vn{%custom(:)} array allocated. [In
this case, none of the custom values have been set so are zero by definition.] In the above example,
the \vn{%custom} array will be allocated if needed in the call to \vn{pointer_to_attribute}.

If not defined through a lattice file, custom attributes can also be
defined directly from within a program using the 
\Hyperref{r:set.custom.attribute.name}{set_custom_attribute_name} routine.
For example:
\begin{example}
  logical err_flag
  ...
  call set_custom_attribute_name ('QUADRUPOLE::ERROR_CURRENT', err_flag) 
\end{example}

Note: When there is a superposition (\sref{c:super}), the \vn{super_slave} elements that are formed
do {\em not} have any custom attributes assigned to them even when their \vn{super_lord} elements
have custom attributes. This is done since the \bmad bookkeeping routines are not able to handle the
situation where a \vn{super_slave} element has multiple \vn{super_lord} elements and thus the custom
attributes from the different \vn{super_lord} elements have to be combined. Proper handling of this
situation is left to any custom code that a program implements to handle custom attributes.

%--------------------------------------------------------------------------
\section{Bmad Reserved Variables}

\index{ele_struct!\%ixx} \index{ele_struct!\%iyy} \index{ele_struct!\%izz} \index{ele_struct!\%bmad_logic}
A number of \vn{ele_struct} components are reserved for \bmad internal use only. These are:
\begin{example}
  %ixx
  %iyy
  %izz
  %bmad_logic
  %select
\end{example}
To avoid conflict with multiple routines trying to use these components simultaneously, these
components are only used for short term bookkeeping within individual routines.

